\section{超越 Prompt：改变模型本身}

\subsection{参数高效微调方法}

当 prompt 优化到达极限时，可以通过修改模型权重来进一步提升性能。但对于拥有数千亿参数的 LLM，完全微调（full fine-tuning）成本极高。因此研究者发展出了一系列\textbf{参数高效方法}（Parameter-Efficient Fine-Tuning, PEFT）。

讲师介绍了几种主流方法：

\begin{table}[H]
\centering
\begin{tabular}{lp{8cm}}
\toprule
\textbf{方法} & \textbf{核心思路} \\
\midrule
Adapter & 在 Transformer 层之间插入小型适配器模块，只训练适配器参数 \\
BitFit  & 仅微调模型中的 bias 参数（数量远少于全部权重） \\
LoRA    & 通过低秩分解（low-rank decomposition）在原权重矩阵旁添加辅助矩阵 $\Delta W = BA$，只训练 $B$ 和 $A$ \\
\bottomrule
\end{tabular}
\caption{主流参数高效微调方法对比}
\end{table}

\begin{importantbox}{LoRA 的架构优势}
LoRA（Low-Rank Adaptation）之所以在工业界广受欢迎，不仅因为训练效率高，更因为它具有优秀的\textbf{架构友好性}：

\begin{itemize}
    \item 原始模型权重 $W$ 保持不变，只需存储小型的辅助矩阵 $\Delta W$
    \item 在推理时，将 $\Delta W$ 投影叠加到 $W$ 上即可：$W' = W + \Delta W$
    \item 一台服务器可运行一个基础模型，同时加载多个不同的 LoRA adapter
    \item 不同 adapter 实现不同功能（如``专业邮件风格'' vs ``莎士比亚十四行诗风格''），按需切换
\end{itemize}

相比之下，如果使用 full fine-tuning，每种功能都需要一个完整的模型副本，资源消耗成倍增长。
\end{importantbox}

\subsection{多种``合法''的语言模型}

讲师引入了一个重要概念：存在多种同样``合法''（valid）的语言模型。

以一个日常例子说明：给定 ``The storage compartment in the back of your car is called a \_\_\_\_''，美式英语用户会说 ``trunk''，而英式英语用户会说 ``boot''。两个答案都是正确的——它们代表了两种不同的、同样合法的语言模型。

类似的例子无处不在：
\begin{itemize}
    \item ``The ruler of the country is called the \_\_\_\_'' $\rightarrow$ king / president / premier
    \item 在美国新泽西州，一种三明治在不同地区被称为 sub / hoagie / hero
    \item 你与朋友、父母、教授说话的方式各不相同——这也是不同的``子语言模型''
\end{itemize}

\begin{knowledgebox}{在合法语言模型之间移动}
在不同``合法语言模型''之间移动（shift）是一个重要的研究课题。它同时服务于两个目的：
\begin{itemize}
    \item \textbf{产品定制}：让模型采用特定的语调和风格（如客服机器人的专业语气）
    \item \textbf{AI Safety}：确保模型拒绝有害请求（从``能够回答任何问题''的语言模型 shift 到``拒绝回答危险问题''的语言模型）
\end{itemize}
无论目的如何，技术手段都是相同的：通过更新权重来改变模型在概率空间中的位置。
\end{knowledgebox}

讲师用 Gemini 做了一个有趣的实验：试图让模型在 ``trunk'' 和 ``boot'' 之间做出选择。结果 Gemini（经过大量 post-training）能够识别出这是一个有歧义的问题，并同时给出两种答案，而不是偏向任何一方。这展示了现代 LLM 经过精心调教后的``多视角''能力。

\subsection{Instruction Tuning（指令微调）}

将语言模型的行为从``复述训练数据''转变为``遵循指令做有用的事''的过程称为 Instruction Tuning。

具体做法是构建一个指令-回答数据集，例如：

\begin{quote}
\textbf{Goal}: Get a cool sleep on a summer day. \\
\textbf{How would you accomplish this goal?} \\
(A) Put a wet towel on the pillow \quad (B) Put a pillow in the oven \\
\textbf{Answer}: (A)
\end{quote}

\begin{importantbox}{Instruction Tuning 的效果}
实验表明，经过 Instruction Tuning 后，模型在多种 held-out（训练时未见过的）任务上性能显著提升，尤其是大模型受益更明显。Instruction Tuning 的本质是教会模型\textbf{以人类认为有用的方式}来组织和呈现其训练数据中的知识，而不仅仅是逐字复述训练数据。
\end{importantbox}

\subsection{RLHF 与 Constitutional AI}

\textbf{RLHF（Reinforcement Learning from Human Feedback）}的流程：
\begin{enumerate}
    \item 让模型对同一 prompt 生成多个回答
    \item 人类标注员对回答进行偏好排序
    \item 训练一个\textbf{奖励模型}（Reward Model）来模拟人类偏好
    \item 用奖励模型作为强化学习的 reward signal，优化语言模型的输出
\end{enumerate}

\textbf{Constitutional AI}（由 Anthropic 提出）则进一步追问：我们是否真的需要人类偏好数据？

\begin{knowledgebox}{Constitutional AI 的核心思想}
Constitutional AI 的做法是：
\begin{enumerate}
    \item 写出一套明确的``宪法''规则（constitution），描述模型应遵守的行为准则
    \item 用一个 LLM 评估另一个 LLM 的输出是否符合这些规则
    \item 根据评估结果更新模型权重
\end{enumerate}
这种方法用 LLM 自身来替代人类标注员的角色，大幅降低了对齐（alignment）的人力成本。讲师引用了 Anthropic 公开的 constitution 内容片段作为示例。
\end{knowledgebox}

讲师总结了这一节的核心信息：无论使用哪种技术——prompt 工程、指令微调、RLHF 还是 Constitutional AI——最终目标都是相同的：\textbf{通过某种方式更新语言模型的权重，使其行为向我们期望的方向偏移}。

\subsection{本章小结}

当 prompt 优化到达极限时，有多种方式可以修改模型本身：LoRA 等参数高效方法适合成本敏感场景；Instruction Tuning 教模型``如何有用地回答''；RLHF 和 Constitutional AI 用偏好信号让模型的输出更符合人类期望。所有这些技术的共同点是通过持续的 next-word prediction 训练来调整权重，在多种``合法语言模型''之间进行选择。

